\documentclass[11pt]{article}

\usepackage[T1]{fontenc}
\usepackage[utf8]{inputenc}
\usepackage{newtxtext,newtxmath}
\usepackage[letterpaper,margin=0.82in]{geometry}
\usepackage{microtype}
\usepackage{amsmath}
\usepackage{booktabs}
\usepackage{enumitem}
\usepackage{xcolor}
\usepackage{hyperref}
\usepackage{fancyhdr}
\usepackage{titlesec}
\usepackage{array}

\definecolor{ink}{HTML}{172033}
\definecolor{accent}{HTML}{315A8A}
\definecolor{soft}{HTML}{EDF3F8}
\definecolor{rulegray}{HTML}{C9D3DE}
\definecolor{cBlue}{HTML}{2563EB}
\definecolor{cGreen}{HTML}{16A34A}
\definecolor{cRed}{HTML}{DC2626}
\hypersetup{
  colorlinks=true,
  linkcolor=accent,
  citecolor=accent,
  urlcolor=accent,
  pdftitle={S4 Permutahedron Continuous-Journey Disassembly},
  pdfauthor={Brandon Kirkland},
  pdfsubject={Implementation of the six-strand disassembly algorithm}
}
\setlength{\parindent}{0pt}
\setlength{\parskip}{0.55em}
\setlist[itemize]{leftmargin=1.35em,itemsep=0.2em,topsep=0.25em}
\setlist[enumerate]{leftmargin=1.5em,itemsep=0.2em,topsep=0.25em}
\titleformat{\section}{\large\bfseries\color{ink}}{}{0pt}{}[\vspace{-0.25em}\color{rulegray}\titlerule]
\titleformat{\subsection}{\normalsize\bfseries\color{accent}}{}{0pt}{}
\pagestyle{fancy}
\fancyhf{}
\lhead{\small $S_4$ permutahedron continuous-journey disassembly}
\rhead{\small Pr\'ecis, August 2026}
\cfoot{\small\thepage}
\renewcommand{\headrulewidth}{0.3pt}

\newcommand{\summarybox}[1]{%
  \begin{center}
  \fcolorbox{rulegray}{soft}{\begin{minipage}{0.94\linewidth}\small #1\end{minipage}}
  \end{center}}
\newcommand{\B}{\textcolor{cBlue}{\mathsf{B}}}
\newcommand{\G}{\textcolor{cGreen}{\mathsf{G}}}
\newcommand{\R}{\textcolor{cRed}{\mathsf{R}}}
\newcommand{\perm}[1]{\mathtt{#1}}

\begin{document}

\begin{center}
  {\LARGE\bfseries\color{ink} $S_4$ Permutahedron Continuous-Journey Disassembly}\par
  \vspace{0.3em}
  {\large Six-strand decomposition via colored journeys on the truncated octahedron}\par
  \vspace{0.55em}
  {\small Pr\'ecis prepared for S.\,J.\ Gates Jr.}
\end{center}

\summarybox{%
Following written instructions dated 2026-08-26, the 24-vertex $S_4$
permutahedron has been decomposed into six four-vertex strands by
tracing an 8-edge continuous colored journey from each vertex of the
base hexagonal face.  Each strand's four recorded vertices form a
right coset of the Klein four-group
$V_4 = \{e,\, (12)(34),\, (14)(23),\, (13)(24)\}$ inside $S_4$.
The six cosets partition all 24 vertices with no overlaps.
An interactive 3D visualization implements the strand-by-strand
separation with animated edge tracing.  Three encoding decisions
require confirmation before the artifact is marked publishable.}

\section{Algorithm}

The procedure, as stated in the written instructions:

\begin{enumerate}
  \item Start at a vertex on the base face and \emph{record} it.
  \item Travel along 2 links using
        \textcolor{cBlue}{blue} and \textcolor{cRed}{red},
        record the vertex reached.
  \item Continue along 4 links in the sequence
        \textcolor{cGreen}{green}, \textcolor{cBlue}{blue},
        \textcolor{cRed}{red}, \textcolor{cGreen}{green},
        record the vertex reached.
  \item Continue along 2 links using
        \textcolor{cBlue}{blue} and \textcolor{cRed}{red},
        record the vertex reached.
  \item The four recorded vertices are one of the six quartets.
  \item Repeat from the next vertex on the base face until all six
        strands are extracted.
\end{enumerate}

\section{Encoding decisions}

\subsection{Color-generator map}

Confirmed from the referenced screenshot:

\begin{center}
\begin{tabular}{@{}lll@{}}
\toprule
Color & Generator & Transposition \\
\midrule
\textcolor{cBlue}{Blue}   & $s_1$ & $(12)$ \\
\textcolor{cGreen}{Green} & $s_2$ & $(23)$ \\
\textcolor{cRed}{Red}     & $s_3$ & $(34)$ \\
\bottomrule
\end{tabular}
\end{center}

\subsection{Base face}

The $s_2,s_3$ hexagonal face where position~1 is fixed at value~1,
confirmed visually from the screenshot (base of the figure):
\[
  \perm{1234} \xrightarrow{s_2} \perm{1324} \xrightarrow{s_3}
  \perm{1342} \xrightarrow{s_2} \perm{1432} \xrightarrow{s_3}
  \perm{1423} \xrightarrow{s_2} \perm{1243} \xrightarrow{s_3}
  \perm{1234}
\]

\subsection{Journey word}

Three segments, eight edges total:

\begin{center}
\begin{tabular}{@{}llcl@{}}
\toprule
Segment & Generators & Length & Net permutation \\
\midrule
$J_1$ & $\B\,\R$ \quad $(s_1,\, s_3)$ & 2 & $(12)(34) = A$ \\
$J_2$ & $\G\,\B\,\R\,\G$ \quad $(s_2,\, s_1,\, s_3,\, s_2)$ & 4 & $(13)(24) = C$ \\
$J_3$ & $\B\,\R$ \quad $(s_1,\, s_3)$ & 2 & $(12)(34) = A$ \\
\bottomrule
\end{tabular}
\end{center}

Full word: $[s_1,\, s_3,\, s_2,\, s_1,\, s_3,\, s_2,\, s_1,\, s_3]$.
Vertices are recorded at cumulative positions $\{0,\, 2,\, 6,\, 8\}$.

\subsection{Klein four-group structure}

The four recorded vertices of each strand form a right coset
$\sigma \cdot V_4$ where:

\begin{center}
\begin{tabular}{@{}lll@{}}
\toprule
Element & Permutation & Double transposition \\
\midrule
$e$  & $\perm{1234}$ & identity \\
$A$  & $\perm{2143}$ & $(12)(34)$ \\
$AC$ & $\perm{4321}$ & $(14)(23)$ \\
$C$  & $\perm{3412}$ & $(13)(24)$ \\
\bottomrule
\end{tabular}
\end{center}

\section{The six strands}

\renewcommand{\arraystretch}{1.35}
\begin{center}
\small
\begin{tabular}{@{}clp{7.2cm}ll@{}}
\toprule
\# & Seed & Route & Quartet & Sector \\
\midrule
1 & $\perm{1234}$ &
  $\perm{1234}\xrightarrow{\B}\perm{2134}\xrightarrow{\R}\perm{2143}\xrightarrow{\G}\perm{2413}\xrightarrow{\B}\perm{4213}\xrightarrow{\R}\perm{4231}\xrightarrow{\G}\perm{4321}\xrightarrow{\B}\perm{3421}\xrightarrow{\R}\perm{3412}$ &
  $P_6$ & VM3 \\
2 & $\perm{1324}$ &
  $\perm{1324}\xrightarrow{\B}\perm{3124}\xrightarrow{\R}\perm{3142}\xrightarrow{\G}\perm{3412}\xrightarrow{\B}\perm{4312}\xrightarrow{\R}\perm{4321}\xrightarrow{\G}\perm{4231}\xrightarrow{\B}\perm{2431}\xrightarrow{\R}\perm{2413}$ &
  $P_3$ & VM \\
3 & $\perm{1342}$ &
  $\perm{1342}\xrightarrow{\B}\perm{3142}\xrightarrow{\R}\perm{3124}\xrightarrow{\G}\perm{3214}\xrightarrow{\B}\perm{2314}\xrightarrow{\R}\perm{2341}\xrightarrow{\G}\perm{2431}\xrightarrow{\B}\perm{4231}\xrightarrow{\R}\perm{4213}$ &
  $P_2$ & TM \\
4 & $\perm{1432}$ &
  $\perm{1432}\xrightarrow{\B}\perm{4132}\xrightarrow{\R}\perm{4123}\xrightarrow{\G}\perm{4213}\xrightarrow{\B}\perm{2413}\xrightarrow{\R}\perm{2431}\xrightarrow{\G}\perm{2341}\xrightarrow{\B}\perm{3241}\xrightarrow{\R}\perm{3214}$ &
  $P_4$ & VM1 \\
5 & $\perm{1423}$ &
  $\perm{1423}\xrightarrow{\B}\perm{4123}\xrightarrow{\R}\perm{4132}\xrightarrow{\G}\perm{4312}\xrightarrow{\B}\perm{3412}\xrightarrow{\R}\perm{3421}\xrightarrow{\G}\perm{3241}\xrightarrow{\B}\perm{2341}\xrightarrow{\R}\perm{2314}$ &
  $P_1$ & CM \\
6 & $\perm{1243}$ &
  $\perm{1243}\xrightarrow{\B}\perm{2143}\xrightarrow{\R}\perm{2134}\xrightarrow{\G}\perm{2314}\xrightarrow{\B}\perm{3214}\xrightarrow{\R}\perm{3241}\xrightarrow{\G}\perm{3421}\xrightarrow{\B}\perm{4321}\xrightarrow{\R}\perm{4312}$ &
  $P_5$ & VM2 \\
\bottomrule
\end{tabular}
\end{center}
\renewcommand{\arraystretch}{1.0}

The six quartets contain all 24 vertices of $S_4$ exactly once.
Each quartet is a right coset $\sigma \cdot V_4$, confirming the
$S_4 / V_4$ quotient structure.

\section{Visualization}

The implementation renders an interactive 3D permutahedron (truncated
octahedron) with the following behavior:

\begin{itemize}
  \item \textbf{Initial state:} all 24 vertices and 36 colored edges
        displayed in the orientation matching the referenced screenshot,
        with the base hexagonal face at the bottom.
  \item \textbf{Strand separation:} a single button click triggers each
        strand extraction.  The 8 journey edges glow sequentially in their
        generator colors
        (\textcolor{cBlue}{blue}, \textcolor{cGreen}{green},
         \textcolor{cRed}{red}),
        pausing at each of the four recorded vertices, then the strand's
        vertices separate outward from the polyhedron.
  \item \textbf{Reassembly:} returns all strands to the assembled state.
  \item \textbf{Overview:} extracts all six strands sequentially with
        a single click.
\end{itemize}

\section{Questions requiring confirmation}

Three encoding decisions are not uniquely determined by the written
instructions and require explicit confirmation before the artifact is
marked publishable.

\subsection{Q1: Blue-Red ordering in outer segments}

The instructions say ``using blue and red links'' for the 2-link outer
segments ($J_1$ and $J_3$) without specifying order.  We implemented
Blue then Red ($s_1$ then $s_3$).

Because $s_1 = (12)$ and $s_3 = (34)$ act on disjoint positions,
they commute: $s_1 s_3 = s_3 s_1 = (12)(34)$.  The endpoint is
identical either way, but the intermediate vertex differs:

\begin{center}
\begin{tabular}{@{}ll@{}}
\toprule
Order & Path from $\perm{1234}$ \\
\midrule
$\B\,\R$ (Blue, Red) & $\perm{1234} \to \perm{2134} \to \perm{2143}$ \\
$\R\,\B$ (Red, Blue) & $\perm{1234} \to \perm{1243} \to \perm{2143}$ \\
\bottomrule
\end{tabular}
\end{center}

\textbf{Is Blue-then-Red the intended traversal order?}

\subsection{Q2: Base face traversal order}

The instructions name $\perm{1234}$, then $\perm{1324}$, then
``the next node on the base.''  We continue around the hexagonal
perimeter:
\[
  \perm{1234} \to \perm{1324} \to \perm{1342} \to
  \perm{1432} \to \perm{1423} \to \perm{1243}
\]

\textbf{Is this the intended sequence of seeds?}

\subsection{Q3: Overall correctness}

\textbf{Are the six quartets listed above correct?}

\end{document}
